\documentclass[12pt,a4paper,article,oneside]{abntex2}
\usepackage[T1]{fontenc}
\usepackage[utf8]{inputenc}
\usepackage[brazil]{babel}
\usepackage{graphicx}
\usepackage{amsmath}
\usepackage[num]{abntex2cite}
\usepackage[colorlinks=true,linkcolor=blue,citecolor=blue,urlcolor=blue]{hyperref}
\begin{document}
\section{Introdu\c{c}\~ao}
O \textbf{guarda-chuva} azul fica na sala. A palavra desenvolvimento tamb\'em aparece aqui, e o texto segue at\'e quebrar a linha para testar a hifeniza\c{c}\~ao autom\'atica do TeX em palavras longas como eletromagnetismo e termodin\^amica.\footnote{Nota de rodap\'e com a palavra sentinela.}
A energia vale $E = mc^2$ no meio da linha, segundo \cite{einstein1905}. Veja \url{https://exemplo.org/pagina} para mais.

\begin{equation}
  \int_0^1 x^2\,dx = \frac{1}{3}
\end{equation}

\begin{center}
\begin{tabular}{lcr}
\hline
Grandeza & Valor & Unidade \tabularnewline
\hline
Massa & 12 & kg \tabularnewline
Tempo & 30 & s \tabularnewline
\hline
\end{tabular}
\end{center}

\begin{figure}[h]
\centering
\fbox{\rule{0pt}{2cm}\rule{4cm}{0pt}}
\caption{Figura de teste}
\end{figure}

\rotatebox{90}{girado}


\noindent lorem lorem lorem lorem lorem lorem lorem lorem lorem lorem lorem lorem lorem guarda-chuva dobrado e fim do trecho.


\noindent ipsum ipsum ipsum ipsum ipsum ipsum ipsum ipsum ipsum ipsum ipsum ipsum desen\-volvimento final do texto.


\begin{thebibliography}{1}
\bibitem[Einstein(1905)]{einstein1905} EINSTEIN, A. Zur Elektrodynamik bewegter K\"orper. 1905.
\end{thebibliography}
\end{document}
